\documentclass[11pt]{article}

\usepackage[letterpaper,margin=0.82in,headheight=15pt]{geometry}
\usepackage[T1]{fontenc}
\usepackage{lmodern}
\usepackage{microtype}
\usepackage{amsmath,amssymb,amsthm,bm,mathtools}
\usepackage{stmaryrd}
\usepackage{booktabs,longtable,tabularx,array,multirow}
\usepackage{enumitem}
\usepackage{xcolor}
\usepackage{fancyhdr}
\usepackage{graphicx}
\usepackage{tikz-cd}
\usepackage{hyperref}
\usepackage[nameinlink,noabbrev]{cleveref}

\definecolor{navy}{HTML}{17324D}
\definecolor{teal}{HTML}{147D92}
\definecolor{orange}{HTML}{D47A27}
\definecolor{softblue}{HTML}{EEF4F7}
\definecolor{softgray}{HTML}{F3F5F6}
\definecolor{darkgray}{HTML}{46545E}
\definecolor{softorange}{HTML}{FFF4E8}
\definecolor{softgreen}{HTML}{EDF7F0}
\definecolor{softred}{HTML}{FCEEEE}

\hypersetup{
  colorlinks=true,
  linkcolor=navy,
  citecolor=teal,
  urlcolor=teal,
  pdftitle={Volume XII: Microscopic Wilson-Line Dynamics, Spectral Cuts, Link-Odd Helicity Matrices, and Second-Order Non-Abelian Convergence},
  pdfauthor={DeuteronWigner project}
}

\pagestyle{fancy}
\fancyhf{}
\fancyhead[L]{\small\color{darkgray}DeuteronWigner formalism - Volume XII}
\fancyhead[R]{\small\color{darkgray}Microscopic Wilson dynamics through order two}
\fancyfoot[L]{\small\color{darkgray}31 July 2026}
\fancyfoot[R]{\small\color{darkgray}\thepage}
\renewcommand{\headrulewidth}{0.4pt}

\setlist[itemize]{leftmargin=1.5em,itemsep=2pt,topsep=3pt}
\setlist[enumerate]{leftmargin=1.7em,itemsep=2pt,topsep=3pt}
\setlength{\parindent}{0pt}
\setlength{\parskip}{5pt}
\setlength{\emergencystretch}{2em}
\renewcommand{\arraystretch}{1.16}
\allowdisplaybreaks

\newtheorem{definition}{Definition}[section]
\newtheorem{axiom}[definition]{Axiom}
\newtheorem{principle}[definition]{Design principle}
\newtheorem{requirement}[definition]{Requirement}
\newtheorem{proposition}[definition]{Proposition}
\newtheorem{theorem}[definition]{Theorem}
\newtheorem{lemma}[definition]{Lemma}
\newtheorem{corollary}[definition]{Corollary}
\newtheorem{remark}[definition]{Remark}
\newtheorem{decision}[definition]{Model decision}

\crefname{definition}{definition}{definitions}
\crefname{axiom}{axiom}{axioms}
\crefname{principle}{design principle}{design principles}
\crefname{requirement}{requirement}{requirements}
\crefname{proposition}{proposition}{propositions}
\crefname{theorem}{theorem}{theorems}
\crefname{lemma}{lemma}{lemmas}
\crefname{corollary}{corollary}{corollaries}
\crefname{remark}{remark}{remarks}
\crefname{decision}{model decision}{model decisions}

\newcommand{\kT}{\bm{k}_{T}}
\newcommand{\pT}{\bm{p}_{T}}
\newcommand{\qT}{\bm{q}_{T}}
\newcommand{\lT}{\bm{\ell}_{T}}
\newcommand{\DT}{\bm{\Delta}_{T}}
\newcommand{\bD}{\bm{b}_{\Delta}}
\newcommand{\bM}{\bm{b}_{\mathrm{TMD}}}
\newcommand{\dd}{\mathrm{d}}
\newcommand{\Tr}{\operatorname{Tr}}
\newcommand{\tr}{\operatorname{tr}}
\newcommand{\diag}{\operatorname{diag}}
\newcommand{\rank}{\operatorname{rank}}
\newcommand{\Span}{\operatorname{span}}
\newcommand{\Ker}{\operatorname{Ker}}
\newcommand{\Ran}{\operatorname{Ran}}
\newcommand{\Cov}{\operatorname{Cov}}
\newcommand{\Var}{\operatorname{Var}}
\newcommand{\Id}{\operatorname{Id}}
\newcommand{\Hom}{\operatorname{Hom}}
\newcommand{\End}{\operatorname{End}}
\newcommand{\Hil}{\mathcal H}
\newcommand{\LF}{\ensuremath{\mathrm{LF}}}
\newcommand{\BLFQ}{\ensuremath{\mathrm{BLFQ}}}
\newcommand{\TMD}{\ensuremath{\mathrm{TMD}}}
\newcommand{\GTMD}{\ensuremath{\mathrm{GTMD}}}
\newcommand{\GPD}{\ensuremath{\mathrm{GPD}}}
\newcommand{\PDF}{\ensuremath{\mathrm{PDF}}}
\newcommand{\QCD}{\ensuremath{\mathrm{QCD}}}
\newcommand{\EFT}{\ensuremath{\mathrm{EFT}}}
\newcommand{\TTN}{\ensuremath{\mathrm{TTN}}}
\newcommand{\TN}{\ensuremath{\mathrm{TN}}}
\newcommand{\WTI}{\ensuremath{\mathrm{WTI}}}
\newcommand{\BRST}{\ensuremath{\mathrm{BRST}}}
\newcommand{\cO}{\mathcal O}
\newcommand{\cM}{\mathcal M}
\newcommand{\cR}{\mathcal R}
\newcommand{\cS}{\mathcal S}
\newcommand{\cT}{\mathcal T}
\newcommand{\cV}{\mathcal V}
\newcommand{\cW}{\mathcal W}
\newcommand{\cE}{\mathcal E}
\newcommand{\cG}{\mathcal G}
\newcommand{\cH}{\mathcal H}
\newcommand{\cK}{\mathcal K}
\newcommand{\cQ}{\mathcal Q}
\newcommand{\cP}{\mathcal P}
\newcommand{\cC}{\mathcal C}
\newcommand{\cZ}{\mathcal Z}
\newcommand{\cD}{\mathcal D}
\newcommand{\cN}{\mathcal N}
\newcommand{\cU}{\mathcal U}
\newcommand{\cF}{\mathcal F}
\newcommand{\cA}{\mathcal A}
\newcommand{\cB}{\mathcal B}
\newcommand{\cL}{\mathcal L}
\newcommand{\cJ}{\mathcal J}
\newcommand{\cI}{\mathcal I}
\newcommand{\eps}{\varepsilon}
\newcommand{\abs}[1]{\left|#1\right|}
\newcommand{\norm}[1]{\left\lVert#1\right\rVert}
\newcommand{\inner}[2]{\left\langle #1\middle|#2\right\rangle}
\newcommand{\avg}[1]{\left\langle #1\right\rangle}
\newcommand{\phys}{\mathrm{phys}}
\newcommand{\eff}{\mathrm{eff}}
\newcommand{\bare}{\mathrm{bare}}
\newcommand{\ind}{\mathrm{ind}}
\newcommand{\ct}{\mathrm{ct}}
\newcommand{\reg}{\mathrm{reg}}
\newcommand{\trunc}{\mathrm{trunc}}
\newcommand{\ren}{\mathrm{ren}}
\newcommand{\num}{\mathrm{num}}
\newcommand{\UV}{\mathrm{UV}}
\newcommand{\IR}{\mathrm{IR}}
\newcommand{\COM}{\mathrm{CM}}
\newcommand{\Lz}{L_z}
\newcommand{\Jz}{J^z}
\newcommand{\PV}{\operatorname{PV}}
\newcommand{\Disc}{\operatorname{Disc}}
\newcommand{\sgn}{\operatorname{sgn}}
\newcommand{\Amp}{\mathbf{Amp}}
\newcommand{\Match}{\mathbf{Match}}
\newcommand{\Red}{\mathbf{Red}}
\newcommand{\Proc}{\mathbf{Proc}}
\newcommand{\Infer}{\mathbf{Infer}}
\newcommand{\Hfour}{\Hil_{qqqg}}
\newcommand{\Hsixg}{\Hil_{qqqgg}}
\newcommand{\Hseven}{\Hil_{qqqggg}}
\newcommand{\Hseaone}{\Hil_{qqqq\bar q g}}
\newcommand{\Hseatwo}{\Hil_{qqqq\bar q gg}}
\newcommand{\StateBundle}{\texttt{H7MicroscopicStateBundle}}
\newcommand{\ParentBundle}{\texttt{MicroscopicGTMDParentBundle}}
\newcommand{\WilsonOperator}{\texttt{OrderResolvedWilsonOperator}}
\newcommand{\SpectralRule}{\texttt{SpectralSupportRule}}
\newcommand{\CutLedger}{\texttt{MicroscopicCutLedger}}
\newcommand{\SoftLedger}{\texttt{SoftRapidityLedger}}
\newcommand{\WilsonManifest}{\texttt{WilsonFockSupportManifest}}
\newcommand{\ConvergenceManifest}{\texttt{WilsonConvergenceManifest}}
\newcommand{\ProvComplex}{\texttt{Provenance2Complex}}
\newcommand{\AssumptionBundle}{\texttt{AssumptionBundle}}
\newcommand{\PredictionPlan}{\texttt{PredictionPlan}}

\newcolumntype{Y}{>{\raggedright\arraybackslash}X}
\newcolumntype{L}[1]{>{\raggedright\arraybackslash}p{#1}}

\title{\vspace{-1.0cm}
{\color{navy}\bfseries Volume XII: Microscopic Wilson-Line Dynamics, Spectral Cuts,\\[2mm]
Link-Odd Helicity Matrices, and Second-Order Non-Abelian Convergence}\\[4mm]
\large Order-resolved gauge transport on the H4/H7 common parent, explicit continuum support,\\
quark--antiquark--gluon color structure, soft-overlap matching, and the nuclear-composition boundary}
\author{DeuteronWigner project}
\date{31 July 2026}

\begin{document}
\maketitle

\begin{center}
\begin{minipage}{0.94\textwidth}
\colorbox{softblue}{\parbox{0.96\textwidth}{
\textbf{Status and purpose.}
Volumes~0--XI established the typed geometric architecture, regulated
light-front Hamiltonian and Fock-space hierarchy, common microscopic quark,
antiquark, and gluon GTMD helicity matrices, and the first-order Wilson-line
and cut machinery.  C12/H5 attached that machinery to the H4 microscopic
parent.  C13/H6 and C14/H7 then supplied the explicit higher Fock sectors and
strict second-order Wilson support needed for quarks, antiquarks, and active
gluons.  C15/N0 is the active implementation package for the first normalized
spin-1 nuclear state.  The present volume is the normative nucleon-side
Wilson specification that C15 must consume without reinterpretation.  It
separates partonic staple transport from nuclear coherent propagation, treats
physical absorptive support through explicit spectral measures, compares
strict Dyson and Magnus representations through second order, and carries
soft/rapidity overlap terms without claiming ultraviolet matching or a
physical TMD scheme.  It does not introduce nuclear rescattering, process
color weights, Collins--Soper evolution, or an all-orders Wilson result.
}}
\end{minipage}
\end{center}

\begin{abstract}
A microscopic TMD model requires more than a complex light-front wave
function or a path label attached to a fitted distribution.  The link-odd
part of a quark, antiquark, or gluon correlator must arise from a declared
Wilson operator, the same gauge interactions that appear in the Hamiltonian,
physical spectral support for an absorptive cut, and the OAM and helicity
interferences selected by the operator projector.  We formulate the complete
order-one and order-two Wilson layer for the DeuteronWigner microscopic
nucleon.  The input is the common H4/H7 nonzero-transfer helicity-matrix
parent.  Future- and past-pointing paths generate their pole prescriptions
from the path and Fourier conventions.  Continuum and finite-volume spectral
rules determine the discontinuity; numerical broadening never constitutes
physical support.  At first order, complete quark and antiquark link-odd
matrices generate distinct Sivers and Boer--Mulders projections, while the
gluon parent retains four ordered link pairs, independent antisymmetric
\(f^{abc}\) and symmetric \(d^{abc}\) color channels, and all transverse
polarization sectors.  At second order, explicit \(qqqgg\), \(qqqggg\), and
\(qqqq\bar q gg\) support permits strict Dyson--Magnus comparisons for all
partonic species.  Ordered single cuts, the real double-cut intersection,
contact and instantaneous terms, and the square-root-soft expansion are
tracked by typed ledgers and count-once two-cells.  We provide the operator
identities, support theorems, convergence obligations, software interfaces,
benchmark suite, negative tests, and the exact contract with the spin-1
nuclear and QCD-matching layers.  The resulting objects remain regulated
finite-basis boundaries, not physical soft-subtracted TMDs.
\end{abstract}

\tableofcontents

\section{Scientific problem, inherited state, and scope}
\label{sec:scope}

\subsection{The transition from a T-even parent to a dynamical link-odd parent}

Volume~XI constructs the common microscopic matrix
\begin{equation}
 \cM^{a,(0)}_{\lambda_a'\Lambda';\lambda_a\Lambda}
 (x,\kT,\DT)
 \label{eq:t-even-parent}
\end{equation}
for \(a=q,\bar q,g\), at zero skewness and Wilson order zero.  Its future and
past path alternatives coincide, so every link-odd projection vanishes.  The
present volume replaces that controlled null sector by the order-resolved
matrix
\begin{equation}
 \cM^{a,[\gamma],\leq n}
 =
 \sum_{r=0}^{n}
 \cM^{a,[\gamma],(r)},
 \qquad n=1,2,
 \label{eq:order-resolved-parent}
\end{equation}
without abandoning the common state, recoil, helicity, color, or projector
identity.

The dynamical chain is
\begin{equation}
 \begin{aligned}
 \{ |\Psi_{H7}\rangle,\cO_a,\gamma,\cR_{\rm spec}\}
 &\longrightarrow
 \{\cU_\gamma^{(1)},\cU_\gamma^{(2)},\CutLedger,\SoftLedger\}
 \\
 &\longrightarrow
 \cM^{a,[\gamma],\leq2}
 \longrightarrow
 \{F_A^{a,[\gamma]}\}.
 \end{aligned}
 \label{eq:full-chain}
\end{equation}
Here \(\cR_{\rm spec}\) is the declared spectral rule.  A path by itself does
not generate an imaginary part, and a complex numerical regulator does not
replace \(\cR_{\rm spec}\).

\begin{decision}[Primary Wilson boundary]
The authoritative microscopic Wilson boundary is the complete target--parton
helicity matrix before named TMD projection.  Scalar Sivers, Boer--Mulders,
and gluonic-pole functions are reductions of that matrix, not independent
state variables.
\end{decision}

\subsection{Inherited microscopic support}

The H7 state contains the ten sectors
\begin{align}
 \Hil_{H7}={}&
 \Hil_{qqq}\oplus\Hil_{qqqg}
 \oplus\Hil_{qqqu\bar u}\oplus\Hil_{qqqd\bar d}
 \oplus\Hil_{qqqgg}
 \oplus\Hil_{qqqu\bar u g}\oplus\Hil_{qqqd\bar d g}
 \nonumber\\
 &\oplus\Hil_{qqqggg}
 \oplus\Hil_{qqqu\bar u gg}
 \oplus\Hil_{qqqd\bar d gg}.
 \label{eq:h7-space}
\end{align}
The state retains exact quark antisymmetry, antiquark anti-fundamental
identity, two- and three-gluon permutation irreducible representations,
complete color-singlet multiplicities, helicity, \(L_z\), \(J^z\), regulator,
and tensor-network identities.

The support table at the beginning of this volume is
\begin{center}
\begin{tabular}{lcc}
\toprule
Active species & Wilson order one & Wilson order two\\
\midrule
quark & explicit \(qqqg\) & explicit \(qqqgg\)\\
antiquark & explicit \(qqqq\bar q g\) & explicit \(qqqq\bar q gg\)\\
gluon & explicit \(qqqgg\) & explicit \(qqqggg\)\\
\bottomrule
\end{tabular}
\end{center}
Wilson order three is unavailable even though a three-gluon sector exists:
the cubic operator, triple-cut, cubic soft, contact, and gauge-completion
system has not been constructed.

\subsection{Explicit nonclaims}

This volume does not claim:
\begin{itemize}
 \item continuum QCD or regulator independence;
 \item an all-orders Wilson line;
 \item a complete Slavnov--Taylor or BRST proof;
 \item a physical soft function or ultraviolet matching coefficient;
 \item a Collins--Soper evolved TMD;
 \item a process-specific hard/color mixture;
 \item nuclear coherent rescattering, shadowing, or tagged final-state
 interactions;
 \item a physical extraction or calibrated prediction.
\end{itemize}

\subsection{Interface with the active nuclear package}

C15/N0 may consume the order-resolved microscopic nucleon parents, but it must
preserve the distinction
\begin{equation}
 \texttt{PARTONIC\_WILSON\_STAPLE}
 \neq
 \texttt{NUCLEAR\_COHERENT\_PROPAGATION}
 \neq
 \texttt{TAGGED\_FSI}.
 \label{eq:parton-nuclear-distinction}
\end{equation}
Only the first object is available from this volume.

\section{Decorated Wilson identity and support theorem}
\label{sec:identity}

\subsection{Complete operator identity}

A microscopic Wilson operator is identified by
\begin{equation}
 \mathfrak w=
 \bigl(
 a,\Gamma,P_{\rm in},P_{\rm out},
 \gamma,R_c,n_W,
 \cR_{\UV},\cR_{\rm rap},\cS_{\rm soft},
 \mu,\zeta,\mathfrak s,
 \mathfrak f,\mathfrak m
 \bigr),
 \label{eq:wilson-id}
\end{equation}
where \(n_W\) is the Wilson order, \(\mathfrak f\) is the Fock-support
certificate, and \(\mathfrak m\) is the microscopic member identity.  For a
gluon operator, \(\gamma\) is an ordered pair \((\gamma_1,\gamma_2)\) and
the identity additionally retains the final \(f\)- or \(d\)-type color
projection when one has been applied.

\begin{requirement}[No nominal identity]
Two Wilson objects may compose only if their species, momentum fibers,
representation, path geometry, path orientation, Wilson order, regulator,
soft/rapidity convention, microscopic member, and Fock-support certificate
agree or are connected by an explicit adapter.
\end{requirement}

\subsection{Support certificate}

For each requested matrix element, define
\begin{equation}
 \mathfrak f=
 \bigl(
 \nu_{\rm in},\nu_{\rm int},\nu_{\rm out},
 n_W,\cV_{\rm attach},\cV_{\rm contact},
 \cR_{\rm spec},\Delta_{\rm rem}
 \bigr),
 \label{eq:support-certificate}
\end{equation}
where \(\nu_{\rm int}\) is the explicit intermediate-sector sequence and
\(\Delta_{\rm rem}\) is any induced-operator remainder.

\begin{theorem}[Wilson-order/Fock-order gate]
A microscopic Wilson contribution of order \(n\) is admissible only when every
intermediate sector and every contact or constrained-field operator required
by the order-\(n\) expansion is either explicit in the state space or supplied
by a matched induced operator with a visible remainder.  Otherwise the result
is \texttt{UNAVAILABLE\_AT\_THIS\_FOCK\_ORDER}.
\end{theorem}

\begin{proof}
The path-ordered product inserts \(n\) gauge-field operators and generates a
sum over intermediate states between successive insertions.  Projecting out a
required intermediate sector changes both the Hamiltonian resolvent and the
operator.  A bare truncation therefore does not represent the same matrix
element.  Feshbach equivalence is restored only after the induced operator and
its norm kernel are included; any remaining orthogonal component is the
reported remainder.
\end{proof}

\subsection{Unavailable is not zero}

A missing support certificate is a statement about the model truncation, not
a dynamical prediction:
\begin{equation}
 \text{support absent}
 \quad\Longrightarrow\quad
 \texttt{UNAVAILABLE},
 \qquad
 \text{not}\quad F=0.
 \label{eq:unavailable-not-zero}
\end{equation}
A zero is physical only after the relevant state and operator blocks exist and
a symmetry, coupling, cut, or interference limit forces the result to vanish.

\section{Oriented paths, pole prescriptions, and antiunitary reversal}
\label{sec:paths}

\subsection{Semi-infinite line convention}

Let
\begin{equation}
 \gamma_\eta(\lambda)=a+\eta\lambda v,
 \qquad \lambda\geq0,
 \qquad \eta=\pm1.
 \label{eq:path-param}
\end{equation}
With Fourier kernel \(e^{-i\ell\cdot x}\), covariant derivative
\(D_\mu=\partial_\mu+igA_\mu\), and gluon momentum flowing into the eikonal
line, the line integral gives
\begin{equation}
 D_\eta(\ell)
 =\frac{1}{v\cdot\ell-i0\eta}
 =\PV\frac{1}{v\cdot\ell}
 +i\eta\pi\delta(v\cdot\ell).
 \label{eq:eikonal-pole}
\end{equation}
The sign is derived from the full convention tuple and is not an independent
process switch.

\subsection{Path composition and inversion}

Bare parallel transport satisfies
\begin{equation}
 U_{\gamma_2\circ\gamma_1}
 =U_{\gamma_2}U_{\gamma_1},
 \qquad
 U_{\gamma^{-1}}=U_\gamma^\dagger.
 \label{eq:path-laws}
\end{equation}
These are geometric identities of the bare connection.  A renormalized GTMD
also carries cusp, rapidity, soft, ultraviolet, and scheme data, so a bare
path inverse is not by itself the full process-reversal map.

\subsection{Antiunitary link reversal}

Let \(\Theta\) denote the full antiunitary adapter.  It acts on
\begin{equation}
 (P',P,\kT,\DT,\Lambda',\Lambda,
 \lambda_a',\lambda_a,\gamma,R_c)
\end{equation}
by complex conjugation, endpoint exchange, transfer and transverse-momentum
reflection, helicity phases, path inversion, and the correct color
representation map.  The common-frame link-even and link-odd matrices are
\begin{align}
 \cM_{\rm even}
 &=\frac12\left[
 \cM^{[+]}+\Theta^{-1}\cM^{[-]}\Theta
 \right],
 \label{eq:even-matrix}\\
 \cM_{\rm odd}
 &=\frac12\left[
 \cM^{[+]}-\Theta^{-1}\cM^{[-]}\Theta
 \right].
 \label{eq:odd-matrix}
\end{align}
A subtraction of arrays that have merely been labeled future and past is not
an implementation of \cref{eq:odd-matrix}.

\begin{proposition}[Exact zero-rescattering limit]
If the gauge coupling, physical cut support, or every required OAM
interference block is disabled, then \(\cM_{\rm odd}=0\) exactly within the
declared finite-basis tolerance.
\end{proposition}

\section{Spectral measures, light-front resolvents, and physical cuts}
\label{sec:spectral}

\subsection{Continuum spectral rule}

For an intermediate family \(X\), define
\begin{equation}
 \cA_\sigma(E_i)
 =
 \int_{E_{\rm th}}^\infty
 \dd E\,
 \frac{\rho_X(E)N_X(E)}{E_i-E+i0\sigma},
 \qquad \sigma=\pm1.
 \label{eq:continuum-amplitude}
\end{equation}
Then
\begin{equation}
 \operatorname{Im}\cA_\sigma(E_i)
 =-\sigma\pi\rho_X(E_i)N_X(E_i)
 \label{eq:continuum-cut}
\end{equation}
inside the support and vanishes below threshold.

\begin{requirement}[Threshold support]
A nonzero absorptive contribution must carry a threshold, spectral density,
intermediate-state identity, and on-shell support condition.  A finite
numerical \(\eps\) is a convergence parameter and may not enter the physical
result identity.
\end{requirement}

\subsection{Finite-volume sequence}

A discretized spectral measure is
\begin{equation}
 \rho_L(E)
 =\sum_n w_{n,L}\,\delta(E-E_{n,L}),
 \label{eq:finite-volume-measure}
\end{equation}
with a declared volume or resolution sequence \(L\to\infty\).  The associated
quadrature must converge to \cref{eq:continuum-cut} for smooth test functions.
Level smearing may be used to approximate the distribution but must vanish
from the content-addressed result identity.

\subsection{Light-front resolvent}

The Hamiltonian representation is
\begin{equation}
 G_\sigma(P_i^-)
 =\frac{1}{P_i^- -H_0+i0\sigma},
 \label{eq:lf-resolvent}
\end{equation}
and
\begin{equation}
 \cA_\sigma
 =\sum_X
 \frac{
 \langle f|V_{\rm eik}|X\rangle
 \langle X|\cO_\Gamma|i\rangle
 }{P_i^- -P_X^-+i0\sigma}.
 \label{eq:lf-amplitude}
\end{equation}
The cut ledger records whether \cref{eq:eikonal-pole} and
\cref{eq:lf-resolvent} represent the same physical on-shell support or two
distinct cuts.

\subsection{Cut equivalence and count-once rule}

Equivalent descriptions are connected by a two-cell
\begin{equation}
 \texttt{EIKONAL\_CUT}
 \Longleftrightarrow
 \texttt{LF\_RESOLVENT\_CUT}
 \quad[\texttt{EQUIVALENT\_COUNT\_ONCE}].
 \label{eq:count-once}
\end{equation}
Distinct ordered intermediate channels remain additive even when their
numerical denominators coincide.

\section{First-order microscopic Wilson action}
\label{sec:first-order}

\subsection{Matrix-level action}

The first-order correction acts on the complete H4 helicity matrix:
\begin{equation}
 \cM^{a,[\eta],(1)}
 =
 \sum_{X,c}
 \int[\dd\Gamma_X]\,
 \cA^{a,c}_{\rm out\leftarrow X}
 D_\eta(\ell)
 \cA^{a,c}_{X\leftarrow\rm in},
 \label{eq:first-order-matrix}
\end{equation}
where the color, helicity, OAM, state-member, cut, and support labels remain
uncontracted until the appropriate projector is applied.

\subsection{Quark and antiquark channels}

For antiquarks,
\begin{equation}
 t_{\bar3}^a=-(t_3^a)^T,
 \qquad
 U_{\bar3}[\gamma]=U_3[\gamma]^*.
 \label{eq:antifundamental}
\end{equation}
The antiquark result is calculated from direct positive-\(x\) H3/H7 active
slots.  It is not copied from a quark matrix or from negative \(x\).

\subsection{Distinct Sivers and Boer--Mulders reductions}

The matrix-level reductions are
\begin{align}
 f_{1T}^{\perp q}
 &\sim
 \Tr\left[P_{\rm Siv}\cM_{\rm odd}^{q}\right],
 \label{eq:sivers-proj}\\
 h_1^{\perp q}
 &\sim
 \Tr\left[P_{\rm BM}\cM_{\rm odd}^{q}\right].
 \label{eq:bm-proj}
\end{align}
The two projectors select different target-spin, active-parton-spin, and OAM
interference blocks.  Their shared rescattering kernel creates correlations,
but no operator identity of the form
\begin{equation}
 h_1^{\perp q}=C_q f_{1T}^{\perp q}
 \label{eq:no-proportionality}
\end{equation}
is permitted.

\subsection{Active-gluon ordered links}

The gluon operator has two ordered adjoint paths,
\begin{equation}
 \mathfrak g=(\gamma_1,\gamma_2,R_A,\eta_1,\eta_2),
 \label{eq:gluon-link-pair}
\end{equation}
with the four classes
\begin{equation}
 [+,+],\quad[+,-],\quad[-,+],\quad[-,-].
 \label{eq:four-link-pairs}
\end{equation}
The common three-adjoint kernel is projected independently onto
\begin{equation}
 K_f=\frac{-if^{abc}K^{abc}}{24},
 \qquad
 K_d=\frac{d^{abc}K^{abc}}{40/3},
 \label{eq:fd-projectors}
\end{equation}
using
\begin{equation}
 f^{abc}f^{abc}=24,
 \qquad
 d^{abc}d^{abc}=\frac{40}{3},
 \qquad
 f^{abc}d^{abc}=0.
 \label{eq:fd-norms}
\end{equation}
Link topology and color contraction are different labels; neither determines
the other.

\section{Strict second-order non-Abelian transport}
\label{sec:second-order}

\subsection{Dyson polynomial}

For a path generator \(\cK_\eta(s)=\mathcal O(g)\), the strict order-two
Dyson operator is
\begin{equation}
 U_{D,\eta}^{[2]}
 =I+
 \int\dd s_1\,\cK_\eta(s_1)
 +
 \int_{s_1>s_2}\dd s_1\dd s_2\,
 \cK_\eta(s_1)\cK_\eta(s_2).
 \label{eq:dyson2}
\end{equation}
No terms of order \(g^3\) or higher are retained.

\subsection{Magnus polynomial}

Define
\begin{align}
 \Omega_1
 &=\int\dd s_1\,\cK_\eta(s_1),
 \label{eq:omega1}\\
 \Omega_2
 &=\frac12
 \int_{s_1>s_2}\dd s_1\dd s_2\,
 [\cK_\eta(s_1),\cK_\eta(s_2)].
 \label{eq:omega2}
\end{align}
The strict order-two Magnus polynomial is
\begin{equation}
 U_{M,\eta}^{[2]}
 =I+\Omega_1+\frac12\Omega_1^2+\Omega_2.
 \label{eq:magnus2}
\end{equation}
This is not the fully exponentiated \(\exp(\Omega_1+\Omega_2)\), which would
contain uncontrolled higher-order terms.

\begin{theorem}[Dyson--Magnus equality through order two]
For the same path generator, ordering convention, representation, and
integration domain,
\begin{equation}
 U_{D,\eta}^{[2]}-U_{M,\eta}^{[2]}=\mathcal O(g^3).
 \label{eq:dyson-magnus-equality}
\end{equation}
\end{theorem}

\begin{proof}
Split \(\Omega_1^2/2\) into the two ordered domains \(s_1>s_2\) and
\(s_2>s_1\).  Adding \(\Omega_2\) cancels the reverse-ordered commutator term
and leaves exactly the path-ordered product in \cref{eq:dyson2}.  Equal-time
sets have measure zero or are represented by separately declared contact
operators.
\end{proof}

\subsection{Why the commutator matters}

For commuting color fields, \(\Omega_2=0\).  For noncommuting fields,
omitting \(\Omega_2\) destroys the order-two equality and changes the
connected color structure.  The commutator benchmark is therefore a mandatory
non-Abelian test, not an optional numerical diagnostic.

\subsection{Composition, reversal, and order-two unitarity}

The strict approximants must satisfy:
\begin{align}
 U^{[2]}(s_3,s_2)U^{[2]}(s_2,s_1)
 &=U^{[2]}(s_3,s_1)+\mathcal O(g^3),
 \label{eq:composition2}\\
 U_{\gamma^{-1}}^{[2]}
 &=U_\gamma^{[2]\dagger}+\mathcal O(g^3),
 \label{eq:reversal2}\\
 U^{[2]\dagger}U^{[2]}
 &=I+\mathcal O(g^3).
 \label{eq:unitarity2}
\end{align}
The defect relative to an exact ordered exponential must scale cubically in a
coupling-rescaling test.

\subsection{Relation to non-Abelian exponentiation}

Non-Abelian exponentiation organizes connected color factors in the logarithm
of products of Wilson lines \cite{Gardi2013}.  It does not justify replacing
the microscopic order-two calculation by an arbitrary exponentiated
one-gluon phase.  The connected-web statement, the Magnus commutator, and the
explicit H7 Fock-support obligations must all be satisfied independently.

\section{Second-order color and link topology}
\label{sec:color2}

\subsection{Fundamental and anti-fundamental products}

The ordered fundamental product decomposes as
\begin{equation}
 T^aT^b
 =\frac{1}{2N_c}\delta^{ab}I
 +\frac12d^{abc}T^c
 +\frac{i}{2}f^{abc}T^c.
 \label{eq:TaTb}
\end{equation}
The reverse order changes only the commutator term.  For antiquarks, every
generator is replaced by the anti-fundamental representation before the
ordered product is formed.

\subsection{Adjoint transport}

The adjoint generators are
\begin{equation}
 (F^a)_{bc}=-if^{abc}.
 \label{eq:adjoint-generators}
\end{equation}
The second-order active-gluon calculation must compare strict Dyson and Magnus
operators in the adjoint representation rather than borrowing the
fundamental result by Casimir rescaling.

\subsection{Two-link insertion topology}

For the gluon correlator, the second-order ledger keeps
\begin{equation}
 LL,
 \qquad RR,
 \qquad LR,
 \qquad RL
 \label{eq:two-link-topologies}
\end{equation}
as distinct insertion classes.  Here the letters refer to the left and right
ordered Wilson links in the field-strength correlator.  The topology is
retained before the final \(f\)- or \(d\)-type contraction.

\subsection{Microscopic color multiplicities}

The state-color multiplicity and the Wilson color class are independent.  The
H7 state retains 22 \(qqqggg\) singlets and 28 singlets in each
\(qqqq\bar qgg\) sector, including symmetric, antisymmetric, and mixed gluon
permutation representations.  These state labels may mix dynamically but may
not be collapsed into the final gluonic-pole \(f/d\) label.

\section{Two-step cuts and unitarity}
\label{sec:cuts2}

\subsection{Distributional decomposition}

For two denominators,
\begin{equation}
 \frac{1}{x+i0\sigma_1}
 \frac{1}{y+i0\sigma_2},
 \label{eq:two-denom}
\end{equation}
the strict distributional product contains
\begin{align}
 &\PV\frac1x\PV\frac1y,
 \label{eq:pvpv}\\
 &-i\pi\sigma_1\delta(x)\PV\frac1y,
 \qquad
 -i\pi\sigma_2\PV\frac1x\delta(y),
 \label{eq:singlecuts}\\
 &-\pi^2\sigma_1\sigma_2\delta(x)\delta(y).
 \label{eq:doublecut}
\end{align}
The double-cut intersection is real.  It is not the square of one delta
distribution.

\subsection{Ordered cut ledger}

Every cut record contains:
\begin{itemize}
 \item the ordered intermediate sectors;
 \item the two energy denominators;
 \item the two support surfaces and their intersection;
 \item the path orientation and representation;
 \item the operator attachments and contact terms;
 \item equivalence relations to eikonal or resolvent descriptions;
 \item the count-once status.
\end{itemize}

\begin{requirement}[No squared delta]
If two nominal denominators represent the same singular surface, the product
must be defined by an explicit finite-volume, regulator, subtraction, or
distribution-extension rule.  The implementation may not form
\(\delta(x)^2\).
\end{requirement}

\subsection{Optical and probability ledgers}

The order-two discontinuity is compared with the sum over declared on-shell
intermediate channels.  This is a finite-model unitarity benchmark, not a
complete continuum optical theorem.  Missing channels remain visible in a
unitarity-remainder ledger.

\section{Soft and rapidity overlap through second order}
\label{sec:soft2}

\subsection{Square-root-soft expansion}

Let
\begin{equation}
 S=1+aS^{(1)}+a^2S^{(2)}+\mathcal O(a^3).
 \label{eq:soft-series}
\end{equation}
Then
\begin{equation}
 S^{-1/2}
 =1-\frac12aS^{(1)}
 +a^2\left[
 \frac38(S^{(1)})^2-\frac12S^{(2)}
 \right]
 +\mathcal O(a^3).
 \label{eq:soft-sqrt}
\end{equation}
For
\(W_{\rm unsub}=W^{(0)}+aW^{(1)}+a^2W^{(2)}+\cdots\),
\begin{align}
 W_{\rm sub}^{(1)}
 &=W^{(1)}-\frac12S^{(1)}W^{(0)},
 \label{eq:sub1}\\
 W_{\rm sub}^{(2)}
 &=W^{(2)}-\frac12S^{(1)}W^{(1)}
 +\left[
 \frac38(S^{(1)})^2-\frac12S^{(2)}
 \right]W^{(0)}.
 \label{eq:sub2}
\end{align}
Ultraviolet and rapidity counterterms are appended according to the declared
scheme and remain separate from the microscopic overlap.

\subsection{Rapidity-cancellation benchmark}

For every fundamental, anti-fundamental, adjoint, and ordered-two-link class,
the analytic benchmark must satisfy
\begin{equation}
 \frac{\partial}{\partial L_{\rm rap}}
 W_{\rm sub}^{(r)}=0,
 \qquad r=1,2,
 \label{eq:rapidity-cancel}
\end{equation}
within the declared truncation.  Missing or duplicated terms must produce
signed residuals.

\subsection{Two admissible routes}

The compiler permits two mutually exclusive choices:
\begin{enumerate}
 \item \texttt{BOUNDARY\_ONLY\_RESCATTERING}: the microscopic calculation
 supplies a low-resolution unsubtracted boundary and the later matching layer
 supplies the complete soft/rapidity scheme;
 \item \texttt{JOINT\_MICROSCOPIC\_SOFT\_SECTOR}: the boundary and soft sector
 are solved and matched together.
\end{enumerate}
Only the first route is presently complete through the analytic order-two
benchmark.  The second remains unavailable.

\subsection{Proper GTMD status}

A physical GTMD requires soft-radiation and rapidity structure beyond the
unsubtracted overlap \cite{Echevarria2016}.  Therefore every output from this
volume retains
\begin{verbatim}
UNSUBTRACTED_OR_ANALYTICALLY_SUBTRACTED_REGULATED
UV_FINITE_MATCHING_REQUIRED
PHYSICAL_TMD_SCHEME_NOT_ASSIGNED
CONTINUUM_SOFT_FUNCTION_INCOMPLETE
NO_COLLINS_SOPER_EVOLUTION
NO_PROCESS_FACTOR_APPLIED
\end{verbatim}

\section{Gauge completion and current consistency}
\label{sec:gauge}

\subsection{Finite gauge ledger}

The order-two finite-basis gauge benchmark closes only when all required
pieces are present:
\begin{equation}
 \delta_{\rm gauge}^{(2)}
 =\delta_{\rm seq}
 +\delta_{3g}
 +\delta_{4g}
 +\delta_{\rm inst,f}
 +\delta_{\rm inst,g}
 +\delta_{\rm pair}
 +\delta_{\chi}
 +\delta_{\rm ct}
 +\delta_J
 +\delta_Z
 +\delta_{\rm reg}.
 \label{eq:gauge-ledger}
\end{equation}
Each nonzero term has a Hamiltonian and operator owner.  The closure may not be
achieved by fitting an independent counterterm to each named TMD.

\subsection{Ward--Takahashi and Slavnov--Taylor status}

The commuting-generator finite-basis benchmark tests a Ward--Takahashi-like
identity.  Non-Abelian three- and four-gluon terms test a larger finite gauge
ledger.  Neither establishes the complete Slavnov--Taylor or BRST system,
which would require the full gauge-fixed continuum operator algebra, ghost
sector where applicable, regulator-restoration identities, and all relevant
Fock sectors.

\subsection{Path derivative and field-strength insertion}

A path deformation generates a field-strength insertion schematically as
\begin{equation}
 \frac{\delta U_\gamma}{\delta\gamma^\mu(s)}
 \sim
 ig\,U_{\gamma_{>s}}F_{\mu\nu}(\gamma(s))
 \dot\gamma^\nu(s)U_{\gamma_{<s}}.
 \label{eq:path-derivative}
\end{equation}
This identity connects staple-dependent OAM and torque observables to the same
microscopic gauge field.  It does not identify a Sivers function with an OAM
moment.

\section{Tensor-network representation and observable convergence}
\label{sec:tn}

\subsection{State and operator networks}

The H7 state is a ten-branch symmetry-adapted TTN.  The Wilson operator is a
separate operator network whose virtual edges retain:
\begin{itemize}
 \item Fock-sector sequence;
 \item color representation and outer multiplicity;
 \item gluon permutation representation;
 \item target and parton helicities;
 \item \(L_z\) and \(J^z\);
 \item path, Wilson order, cut, and soft identities.
\end{itemize}
The operator may not be absorbed into a redefinition of the state tensor.

\subsection{Full-bond closure}

For every supported species and Wilson order,
\begin{equation}
 \cM_{\rm exact}^{(r)}
 =\cM_{\rm full\ bond}^{(r)}
 \label{eq:full-bond}
\end{equation}
within the declared tolerance.  This equality is checked before and after
link-even/link-odd projection, color projection, and soft subtraction.

\subsection{Observable-sensitive bond convergence}

Reduced-bond studies report separately:
\begin{equation}
 \left\{
 \Delta E,
 1-|\langle\Psi_\chi|\Psi\rangle|^2,
 \Delta\norm{\cM_{\rm odd}},
 \Delta f_{1T}^{\perp},
 \Delta h_1^\perp,
 \Delta K_f,
 \Delta K_d,
 \Delta_{\rm cut},
 \Delta_{\rm soft}
 \right\}.
 \label{eq:bond-observables}
\end{equation}
A state may be converged for the energy and unconverged for a Wilson-sensitive
OAM interference.  No energy-only acceptance is allowed.

\subsection{Topology and recoupling}

Alternative TTN fusion trees are distinct approximation branches at finite
bond.  They are compared through explicit recoupling maps retaining all color
and permutation multiplicities.  A topology difference is not included as a
statistical parameter uncertainty; it is a controlled model/truncation axis.

\section{Common link-odd helicity matrices and reductions}
\label{sec:projections}

\subsection{Matrix before scalar functions}

For each species and order, the stored object is
\begin{equation}
 \cM^{a,(r)}_{\lambda_a'\Lambda';\lambda_a\Lambda}
 (x,\kT,\DT;\gamma,\mathfrak f),
 \label{eq:stored-matrix}
\end{equation}
with the full target--parton helicity structure.  The scalar GTMD or TMD
registry is a view generated by the same Gram-projector engine used at Wilson
order zero.

\subsection{Generic and degenerate kinematics}

At generic kinematics, the complete quark and antiquark matrix is projected
onto the sixteen-function basis.  At collinear or zero-transfer points, the
projector rank drops and only identifiable combinations are returned.  Link-
odd dynamics does not restore coefficients that are kinematically
unidentifiable.

\subsection{Forward TMD limit}

The regulated TMD boundary is
\begin{equation}
 \Phi^{a,[\gamma],(r)}(x,\kT)
 =\cM^{a,[\gamma],(r)}(x,\kT,\DT=0).
 \label{eq:tmd-limit-wilson}
\end{equation}
This is still a regulated microscopic boundary.  The proper soft-subtracted
QCD TMD is produced only after the Volume~V matching and evolution contract is
satisfied.

\subsection{Twist-three moment adapters}

The first transverse moments of link-odd TMDs may be matched to quark--gluon or
tri-gluon twist-three operators.  The adapter must retain operator
normalization, Wilson path, color class, ultraviolet scheme, and mixing.  A
bare cutoff moment is not equated directly with a renormalized Qiu--Sterman or
tri-gluon matrix element.

\subsection{OAM and torque}

Straight- and staple-link OAM definitions differ by a gauge-field torque term.
The volume therefore carries
\begin{equation}
 L_z^{\rm staple}-L_z^{\rm straight}
 =\cT_z^{\rm torque}[\gamma]
 \label{eq:oam-torque}
\end{equation}
as a matched operator relation with its own support and scheme status.  The
same OAM blocks can feed both torque and link-odd projections without making
the observables identical.

\section{Provenance two-complex and assumption compiler}
\label{sec:provenance}

\subsection{Cells}

The Wilson provenance complex contains:
\begin{itemize}
 \item 0-cells: states, Fock sectors, Wilson operators, cuts, soft factors,
 schemes, and results;
 \item 1-cells: insertion, projection, matching, Feshbach elimination,
 subtraction, and nuclear handoff maps;
 \item 2-cells: explicit/induced equivalence, cut count-once, soft-overlap
 subtraction, Dyson--Magnus equality, and branch alternatives.
\end{itemize}

\subsection{Required two-cells}

The minimum order-two relations are
\begin{align}
 \text{explicit higher Fock sector}
 &\Longleftrightarrow
 \text{induced operator}+\text{remainder},
 \label{eq:explicit-induced-cell}\\
 \text{eikonal cut}
 &\Longleftrightarrow
 \text{resolvent cut}
 \quad[\text{count once}],
 \label{eq:cut-cell}\\
 W_{\rm unsub}^{(r)}
 &\longrightarrow
 W_{\rm sub}^{(r)}
 \quad[\text{soft overlap removed}],
 \label{eq:soft-cell}\\
 U_D^{[2]}
 &\Longleftrightarrow
 U_M^{[2]}
 \quad[\mathcal O(g^3)\text{ remainder}].
 \label{eq:dyson-magnus-cell}
\end{align}

\subsection{Assumption axes}

An \AssumptionBundle{} chooses:
\begin{itemize}
 \item the H7 Hamiltonian plan and microscopic member;
 \item Wilson order and representation;
 \item path geometry and link pair;
 \item spectral-support route;
 \item explicit versus induced support where both exist;
 \item soft-overlap route;
 \item exact or TTN state representation and bond policy.
\end{itemize}
Mutually exclusive branches are compared as complete theories.  They are not
added together.

\section{Contract with spin-1 nuclear composition}
\label{sec:nuclear-contract}

\subsection{Nucleon object delivered to C15/N0}

The nuclear layer receives an order-resolved parent
\begin{equation}
 \begin{aligned}
 \cW_{a/N}^{(0\ldots2)}=\bigl(&
 \cM^{a,(0)},\cM^{a,(1)},\cM^{a,(2)},\mathfrak w,\mathfrak f,
 \\
 &\CutLedger,\SoftLedger,\ConvergenceManifest\bigr).
 \end{aligned}
 \label{eq:nucleon-handoff}
\end{equation}
It retains correlated proton/neutron member identity and may not be reduced to
an independently normalized named TMD before the nuclear spin convolution.

\subsection{Nuclear composition preserves link identity}

For the nucleonic deuteron sector,
\begin{equation}
 \begin{aligned}
 W_{a/D}^{(r)}
 ={}&\sum_{N=p,n}\sum_{\lambda',\lambda}
 \int\frac{\dd y}{y}\frac{\dd^2\pT}{(2\pi)^2}
 \rho^{N/D}_{\Lambda'\Lambda;\lambda'\lambda}
 (y,\pT,\DT)
 \\
 &\times
 W_{a/N,\lambda'\lambda}^{(r)}
 \left(\frac{x_D}{y},
 \kT-\frac{x_D}{y}\pT,
 \DT\right)J_D.
 \end{aligned}
 \label{eq:nuclear-wilson-convolution}
\end{equation}
The nuclear map must retain future/past orientation, ordered gluon link pairs,
\(f/d\) color class, Wilson order, and support status.

\subsection{No nuclear phase by scalar dressing}

C15/N0 may not create a deuteron Sivers, Boer--Mulders, or gluonic-pole
function by multiplying the nucleon result by an arbitrary scalar nuclear
response.  Spin recoupling acts on the complete helicity matrix.

\subsection{Unavailable nuclear mechanisms}

The following remain unavailable in N0:
\begin{verbatim}
NUCLEAR_COHERENT_PROPAGATION
NUCLEAR_GLAUBER_RESCATTERING
TAGGED_FINAL_STATE_INTERACTION
POLARIZED_OR_TENSOR_SHADOWING
NNPI_WILSON_INTERFERENCE
\end{verbatim}
Their absence cannot be compensated by modifying the partonic staple phase.

\section{Contract with LF-to-QCD matching, evolution, and processes}
\label{sec:qcd-contract}

\subsection{Matching map}

The regulated result enters
\begin{equation}
 \widetilde W_{a/h}^{\QCD}
 =Z_{a}^{\LF\to\QCD}
 \otimes_x
 \widetilde W_{a/h}^{\LF,\rm sub}
 +\Delta_{a/h}^{\trunc}.
 \label{eq:lf-qcd-match}
\end{equation}
The matching basis is closed under quark--gluon and twist-three mixing.  It
carries Wilson order, rank, link class, color class, and target polarization.

\subsection{Evolution preserves matrix identities}

Collins--Soper and \(\mu\)-evolution act on the matched operator, not on an
unsubtracted C14 boundary.  Evolution must preserve:
\begin{itemize}
 \item exact future/past reversal;
 \item transverse rank;
 \item quark/antiquark separation;
 \item gluon \(f/d\) channel identity;
 \item correlated microscopic member identity.
\end{itemize}

\subsection{Process map}

A process record supplies the hard channel, measurement, link geometry, color
weights, soft prescription, Glauber status, and factorization status.  Until
such a record exists, there is no process-level mixture
\begin{equation}
 F_g^{\rm process}
 =C_f^{\rm process}F_g^{(f)}
 +C_d^{\rm process}F_g^{(d)}.
 \label{eq:process-mixture}
\end{equation}

\subsection{Factorization status}

The allowed process labels remain
\begin{verbatim}
ESTABLISHED
CONDITIONAL
EXPLORATORY
BROKEN
UNASSESSED
\end{verbatim}
A process marked \texttt{BROKEN} or \texttt{UNASSESSED} cannot consume a
universal separate-hadron TMD formula.

\section{Formal software interfaces}
\label{sec:software}

\begin{longtable}{L{0.32\textwidth}L{0.58\textwidth}}
\toprule
Object & Responsibility\\
\midrule
\endhead
\texttt{OrientedWilsonPath} & Path geometry, endpoints, orientation,
representation, cusp and closure identity.\\
\texttt{PolePrescription} & Convention-derived principal-value and delta
sign; never a caller-selected sign.\\
\SpectralRule & Threshold, continuum or finite-volume measure, quadrature,
and below-threshold status.\\
\texttt{LightFrontResolvent} & Hamiltonian denominator, intermediate-state
identity, regulator, and spectral attachment.\\
\CutLedger & Ordered cut surfaces, intersections, equivalence two-cells, and
count-once rules.\\
\WilsonOperator & Strict Dyson/Magnus order, representation, link topology,
Fock support, contact terms, and microscopic member.\\
\texttt{LinkParityAdapter} & Complete antiunitary transformation into a common
kinematic frame.\\
\texttt{LinkOddHelicityMatrix} & Full target--parton matrix prior to scalar
projection.\\
\texttt{OrderedGluonLinkPair} & Left/right adjoint paths and insertion
topology.\\
\texttt{ColorProjection} & Independent \(f\), \(d\), singlet, commutator,
anticommutator, and orthogonal residual sectors.\\
\SoftLedger & Order-resolved square-root-soft, rapidity, ultraviolet, and
scheme terms.\\
\WilsonManifest & Explicit/induced/unavailable support and remainders.\\
\ConvergenceManifest & Exact/TTN, bond, spectral, cut, Wilson-order, soft,
and gauge residuals.\\
\ProvComplex & State/operator/cut/soft/matching equivalences and exclusions.\\
\bottomrule
\end{longtable}

\section{Stable formal requirements}
\label{sec:requirements}

The implementation should expose stable identifiers.  The minimum normative
set is:
\begin{longtable}{L{0.17\textwidth}L{0.75\textwidth}}
\toprule
ID & Requirement\\
\midrule
\endhead
V12.001 & Every Wilson result consumes a complete H4/H7 helicity parent.\\
V12.002 & Path orientation derives the pole prescription from the convention tuple.\\
V12.003 & Numerical \(\eps\) is absent from physical result identity.\\
V12.004 & Every nonzero imaginary part has declared spectral support.\\
V12.005 & Below-threshold absorption vanishes exactly.\\
V12.006 & Finite-volume spectral sequences converge to an analytic oracle where available.\\
V12.007 & Antiunitary reversal transforms endpoints, momenta, helicities, path, and color.\\
V12.008 & Link-even and link-odd matrices are formed before scalar projection.\\
V12.009 & Quark and antiquark active slots remain distinct.\\
V12.010 & Sivers and Boer--Mulders projectors remain independent.\\
V12.011 & Four ordered gluon link pairs remain distinct.\\
V12.012 & Gluon \(f\)- and \(d\)-type channels remain independent.\\
V12.013 & State color multiplicity is distinct from Wilson color class.\\
V12.014 & Strict Dyson and strict Magnus polynomials agree through order two.\\
V12.015 & The Magnus commutator is required for noncommuting fields.\\
V12.016 & Order-two composition, reversal, and unitarity defects scale as \(g^3\).\\
V12.017 & Quark, antiquark, and gluon order-two support is explicit.\\
V12.018 & Wilson order three fails closed.\\
V12.019 & Every omitted Fock sector has an induced-operator remainder or is unavailable.\\
V12.020 & Single-cut surfaces and the double-cut intersection are stored separately.\\
V12.021 & Squared delta distributions are prohibited.\\
V12.022 & Equivalent eikonal/resolvent cuts are counted once.\\
V12.023 & Physically distinct ordered cuts remain additive.\\
V12.024 & First-order square-root-soft subtraction uses the same operator identity.\\
V12.025 & Second-order subtraction includes \(S^{(1)}W^{(1)}\), \((S^{(1)})^2W^{(0)}\), and \(S^{(2)}W^{(0)}\).\\
V12.026 & Missing and duplicate soft terms produce signed residuals.\\
V12.027 & UV finite matching remains explicitly unresolved.\\
V12.028 & Gauge closure contains propagating, instantaneous, contact, current, and counterterm pieces.\\
V12.029 & Finite gauge closure is not labeled full Slavnov--Taylor or BRST closure.\\
V12.030 & Full-bond TTN reproduces exact Wilson matrices.\\
V12.031 & Bond convergence is assessed on Wilson observables, not only energy.\\
V12.032 & Projector rank loss is reported honestly at degenerate kinematics.\\
V12.033 & Twist-three moment adapters retain matching and mixing status.\\
V12.034 & Explicit and induced descriptions are alternatives with visible remainders.\\
V12.035 & Partonic Wilson transport is distinct from nuclear coherent propagation.\\
V12.036 & Nuclear composition preserves path, order, link pair, and color class.\\
V12.037 & No scalar nuclear response generates a link-odd function.\\
V12.038 & Unmatched Wilson objects cannot enter Collins--Soper evolution.\\
V12.039 & Process color weights require a qualified process map.\\
V12.040 & Every output retains the microscopic member and assumption-plan identity.\\
V12.041 & All manifests are deterministic and content-addressed.\\
V12.042 & Production artifacts and the canonical 216-route registry remain immutable.\\
\bottomrule
\end{longtable}

\section{Benchmark and falsification program}
\label{sec:benchmarks}

\subsection{W-A: oriented semi-infinite line}
Verify \cref{eq:eikonal-pole}, path inversion, principal-value equality, and
opposite future/past cuts.

\subsection{W-B: analytic continuum threshold}
Compare the distributional cut with an analytic spectral density, including
exact below-threshold zero.

\subsection{W-C: finite-volume spectral convergence}
Demonstrate monotone or controlled convergence of the discretized measure and
separate volume, quadrature, and smearing errors.

\subsection{W-D: antiunitary matrix reversal}
Apply the full adapter to generic complex helicity matrices; incomplete
endpoint, momentum, helicity, or color transformations must fail.

\subsection{W-E: quark Sivers/Boer--Mulders distinction}
Use one microscopic kernel and two projectors.  Disable each required spin/OAM
block separately and verify the correct null pattern.

\subsection{W-F: direct antiquark channel}
Verify anti-fundamental color, positive-\(x\) support, charge-conjugation
signatures, and absence of quark copying.

\subsection{W-G: active-gluon ordered links}
Reconstruct all four link pairs, independent \(f/d\) channels, and the trace,
helicity, and symmetric-traceless polarization sectors.

\subsection{W-H: strict Dyson--Magnus equality}
Use commuting and noncommuting piecewise-constant paths in fundamental,
anti-fundamental, adjoint, and ordered-two-link representations.

\subsection{W-I: commutator ablation}
Delete \(\Omega_2\); the noncommuting benchmark must produce a nonzero defect.

\subsection{W-J: order-two cut surfaces}
Separate both single cuts and the real double-cut intersection without a
squared delta.

\subsection{W-K: first- and second-order soft cancellation}
Test the complete square-root-soft series and all missing/duplicate-term
ablations.

\subsection{W-L: gauge-completion ledger}
Remove each sequential, three-/four-gluon, instantaneous, pair, chiral,
counterterm, current, residue, and regulator contribution in turn.

\subsection{W-M: Wilson/Fock support}
Verify explicit quark, antiquark, and gluon support through order two and exact
failure at order three.

\subsection{W-N: exact/full-bond/reduced-bond convergence}
Demonstrate that a reduced bond can retain a small energy error while losing a
large link-odd or commutator-sensitive signal.

\subsection{W-O: nuclear handoff}
Compose the nucleon parent with a spin-resolved two-body spectral matrix while
preserving all path and color identities; reject nuclear coherent phase
requests.

\subsection{W-P: matching and process gates}
Reject evolution, physical-TMD, process-mixture, and production-promotion
requests until their independent contracts are satisfied.

\section{Negative-test catalogue}
\label{sec:negative}

The implementation should include at least 160 ordered failures spanning:
\begin{enumerate}
 \item missing path orientation, endpoint, Fourier, coupling, or momentum-flow data;
 \item caller-supplied pole signs inconsistent with the path;
 \item physicalized numerical \(\eps\);
 \item below-threshold nonzero absorption;
 \item finite-volume sequences without a continuum target or error estimate;
 \item incomplete antiunitary reversal;
 \item quark/antiquark aliasing;
 \item Sivers/Boer--Mulders projector aliasing;
 \item lost ordered gluon link identity;
 \item implicit \(f+d\) mixtures;
 \item incorrect \(f/d\) normalization or nonzero cross contraction;
 \item incomplete state-color multiplicities;
 \item Wilson-order/Fock-order mismatch;
 \item induced support without a transformed operator and remainder;
 \item strict Dyson compared with a fully exponentiated Magnus result;
 \item omitted Magnus commutator for noncommuting fields;
 \item order leakage into an order-two result;
 \item bad path composition, reversal, or unitarity scaling;
 \item duplicate or missing cuts;
 \item squared delta distributions;
 \item missing or duplicated first-/second-order soft terms;
 \item soft and collinear objects with mismatched schemes;
 \item gauge closure with omitted instantaneous/contact/current pieces;
 \item false Slavnov--Taylor or BRST readiness;
 \item reduced-bond acceptance based only on energy;
 \item pseudoinverse presented as full GTMD determination;
 \item direct bare moment labeled a renormalized twist-three function;
 \item partonic Wilson phase used as nuclear shadowing or tagged FSI;
 \item nuclear scalar multiplication used to create link-odd TMDs;
 \item unmatched boundary sent to evolution;
 \item process color weights without a process record;
 \item production registry or authoritative-artifact mutation.
\end{enumerate}

\section{Staged implementation and release sequence}
\label{sec:stages}

\subsection{H5.0: microscopic spectral and cut layer}
Attach continuum and finite-volume spectral rules to the H4 parent and verify
matrix-level first-order cuts.

\subsection{H5.1: all-species first-order link-odd matrices}
Implement quark, antiquark, and gluon matrices and distinct scalar
projections.

\subsection{H5.2: explicit higher-Fock support}
Replace induced antiquark and gluon first-order channels by explicit sectors
with visible comparison remainders.

\subsection{H6.0: strict quark order-two benchmark}
Construct \(qqqgg\) support, Dyson/Magnus equality, two-step cuts, and
second-order soft subtraction.

\subsection{H7.0: all-species order-two closure}
Add \(qqqggg\) and \(qqqq\bar qgg\) sectors and validate antiquark and gluon
order-two transport.

\subsection{N0 handoff}
Freeze the order-resolved nucleon bundle and allow the spin-1 nuclear layer to
consume it without introducing nuclear coherent dynamics.

\section{Formal acceptance criteria}
\label{sec:acceptance}

The Volume XII layer is complete at its declared scope only when:
\begin{enumerate}
 \item one microscopic H7 state supplies all quark, antiquark, and gluon
 helicity parents;
 \item path orientation derives every eikonal pole sign;
 \item nonzero absorption has declared continuum or finite-volume support;
 \item finite numerical broadening is absent from physical identity;
 \item future/past reversal closes at the complete matrix level;
 \item Sivers and Boer--Mulders remain distinct projections of one kernel;
 \item all four gluon link pairs and independent \(f/d\) channels close;
 \item all first- and second-order species channels have explicit Fock support;
 \item strict Dyson and Magnus representations agree through order two;
 \item the commutator ablation is nonzero for noncommuting paths;
 \item order-two composition, reversal, unitarity, and exact-exponential defects
 scale at cubic order;
 \item both single cuts and the real double-cut intersection are represented
 without squared deltas;
 \item equivalent cuts are counted once and distinct channels remain additive;
 \item first- and second-order soft-overlap cancellations close;
 \item UV matching and the physical TMD scheme remain explicitly unresolved;
 \item the finite gauge ledger closes only with all declared pieces;
 \item exact and full-bond TTN Wilson matrices agree;
 \item reduced-bond errors are reported for Wilson observables;
 \item Wilson order three fails closed;
 \item C15 can consume the nucleon parent while nuclear coherent mechanisms
 remain unavailable;
 \item LF-to-QCD matching, evolution, processes, inference, and production
 remain fail-closed;
 \item all manifests are deterministic and all prior authoritative artifacts
 remain byte-identical.
\end{enumerate}

\section{Recommended decision and scientific status}
\label{sec:decision}

The adopted microscopic Wilson layer is an
\emph{order-resolved, support-certified, matrix-level Wilson calculation on
the common H4/H7 light-front state, with convention-derived poles, physical
spectral cuts, independent quark/antiquark/gluon color structure, strict
Dyson--Magnus agreement through second order, typed cut and soft ledgers, and
observable-sensitive tensor-network convergence}.

This establishes a controlled finite-basis microscopic source for link-odd
GTMD and TMD boundaries.  It does not yet establish a physical TMD.  The next
scientific step is nuclear composition at N0, followed by explicit nonnucleonic
and coherent nuclear sectors, and then LF-to-QCD matching, evolution, process
factorization, and shared inference.

\appendix

\section{Distribution and sign dictionary}
\label{app:signs}

The project convention is
\begin{equation}
 \frac{1}{x\pm i0}
 =\PV\frac1x\mp i\pi\delta(x).
\end{equation}
For the path convention of \cref{eq:path-param},
\begin{equation}
 \frac{1}{v\cdot\ell-i0\eta}
 =\PV\frac{1}{v\cdot\ell}
 +i\eta\pi\delta(v\cdot\ell).
\end{equation}
A change of Fourier or momentum-flow convention must be represented by an
explicit adapter that changes both the stored convention identity and every
derived sign.

\section{Color dictionary}
\label{app:color}

For \(SU(N_c)\),
\begin{align}
 [T^a,T^b]&=if^{abc}T^c,\\
 \{T^a,T^b\}&=\frac1{N_c}\delta^{ab}I+d^{abc}T^c,\\
 \Tr(T^aT^b)&=\frac12\delta^{ab}.
\end{align}
For \(SU(3)\),
\begin{equation}
 C_F=\frac43,
 \qquad C_A=3,
 \qquad f^{abc}f^{abc}=24,
 \qquad d^{abc}d^{abc}=\frac{40}{3}.
\end{equation}

\section{Strict Dyson--Magnus audit}
\label{app:dyson-magnus}

The equality through order two follows from
\begin{align}
 \frac12\Omega_1^2
 &=\frac12\int_{s_1>s_2}\dd s_1\dd s_2
 \left[
 \cK_1\cK_2+\cK_2\cK_1
 \right],\\
 \Omega_2
 &=\frac12\int_{s_1>s_2}\dd s_1\dd s_2
 \left[
 \cK_1\cK_2-\cK_2\cK_1
 \right],
\end{align}
so
\begin{equation}
 \frac12\Omega_1^2+\Omega_2
 =\int_{s_1>s_2}\dd s_1\dd s_2\,\cK_1\cK_2.
\end{equation}
The implementation stores the discarded order explicitly and never compares
objects with different perturbative content.

\section{Wilson support schema}
\label{app:support-schema}

\begin{verbatim}
{
  "support_manifest_id": "...",
  "microscopic_state_id": "...",
  "species": "Q | QBAR | G",
  "wilson_order": 0,
  "path_or_link_pair_id": "...",
  "representation": "FUNDAMENTAL | ANTIFUNDAMENTAL | ADJOINT",
  "initial_sector": "...",
  "intermediate_sector_sequence": ["..."],
  "final_sector": "...",
  "contact_operator_ids": ["..."],
  "support_status": "EXPLICIT_FOCK_SUPPORTED | INDUCED_WITH_REMAINDER | UNAVAILABLE",
  "remainder_norm": 0.0,
  "regulator_id": "...",
  "assumption_plan_id": "..."
}
\end{verbatim}

\section{Spectral and cut schema}
\label{app:cut-schema}

\begin{verbatim}
{
  "cut_ledger_id": "...",
  "spectral_rule_id": "...",
  "threshold": 0.0,
  "continuum_density_id": "...",
  "finite_volume_sequence": ["..."],
  "ordered_denominators": ["..."],
  "single_cut_surfaces": ["..."],
  "double_cut_intersection": "...",
  "equivalence_two_cells": ["..."],
  "count_once_certificate": "...",
  "below_threshold_residual": 0.0,
  "continuum_convergence_residual": 0.0,
  "epsilon_is_physical": false
}
\end{verbatim}

\section{Soft and rapidity schema}
\label{app:soft-schema}

\begin{verbatim}
{
  "soft_ledger_id": "...",
  "operator_id": "...",
  "wilson_order": 2,
  "route": "BOUNDARY_ONLY_RESCATTERING",
  "W0_id": "...",
  "W1_id": "...",
  "W2_id": "...",
  "S1_id": "...",
  "S2_id": "...",
  "rapidity_counterterm_ids": ["..."],
  "uv_counterterm_ids": ["..."],
  "rapidity_derivative_residual": 0.0,
  "uv_finite_matching_status": "UNRESOLVED_NOT_ZERO",
  "physical_tmd_scheme_status": "NOT_ASSIGNED"
}
\end{verbatim}

\section{Nuclear handoff schema}
\label{app:nuclear-schema}

\begin{verbatim}
{
  "nucleon_wilson_bundle_id": "...",
  "proton_member_id": "...",
  "neutron_member_id": "...",
  "correlation_member_id": "...",
  "supported_orders": [0, 1, 2],
  "species": ["u", "d", "ubar", "dbar", "g"],
  "path_ids": ["..."],
  "gluon_link_pair_ids": ["..."],
  "gluon_color_classes": ["F_TYPE", "D_TYPE"],
  "partonic_wilson_only": true,
  "nuclear_coherent_propagation": "UNAVAILABLE",
  "tagged_fsi": "UNAVAILABLE",
  "matching_status": "REGULATED_FINITE_BASIS"
}
\end{verbatim}

\section{Compact acceptance checklist}
\label{app:checklist}

Before release, verify:
\begin{enumerate}
 \item complete matrix parent, not scalar-only output;
 \item convention-derived pole signs;
 \item explicit spectral support and exact below-threshold zero;
 \item full antiunitary future/past transformation;
 \item distinct Sivers and Boer--Mulders projectors;
 \item direct positive-\(x\) antiquarks;
 \item four gluon link pairs and independent \(f/d\) channels;
 \item explicit Fock support through order two;
 \item strict Dyson--Magnus equality and nonzero commutator ablation;
 \item cut count-once and no squared delta;
 \item complete first-/second-order soft subtraction;
 \item finite gauge ledger with qualified status;
 \item exact/full-bond equality and observable-sensitive bond study;
 \item order three unavailable;
 \item partonic/nuclear rescattering separation;
 \item matching, evolution, process, inference, and production gates closed.
\end{enumerate}

\small
\begin{thebibliography}{99}
\raggedright

\bibitem{Collins2002}
J.~C.~Collins,
``Leading-twist single-transverse-spin asymmetries: Drell--Yan and
deep-inelastic scattering,''
Phys. Lett. B \textbf{536}, 43 (2002),
arXiv:hep-ph/0204004.

\bibitem{BoerMuldersPijlman2003}
D.~Boer, P.~J.~Mulders, and F.~Pijlman,
``Universality of T-odd effects in single spin and azimuthal asymmetries,''
Nucl. Phys. B \textbf{667}, 201 (2003),
arXiv:hep-ph/0303034.

\bibitem{Bomhof2006}
C.~J.~Bomhof, P.~J.~Mulders, and F.~Pijlman,
``The construction of gauge-links in arbitrary hard processes,''
Eur. Phys. J. C \textbf{47}, 147 (2006),
arXiv:hep-ph/0601171.

\bibitem{Meissner2009}
S.~Mei{\ss}ner, A.~Metz, and M.~Schlegel,
``Generalized parton correlation functions for a spin-1/2 hadron,''
JHEP \textbf{08}, 056 (2009),
arXiv:0906.5323.

\bibitem{LorcePasquini2013}
C.~Lorc\'e and B.~Pasquini,
``Structure analysis of the generalized correlator of quark and gluon for a
spin-1/2 target,''
JHEP \textbf{09}, 138 (2013),
arXiv:1307.4497.

\bibitem{Echevarria2016}
M.~G.~Echevarria, A.~Idilbi, K.~Kanazawa, C.~Lorc\'e, A.~Metz,
B.~Pasquini, and M.~Schlegel,
``Proper definition and evolution of generalized transverse momentum
distributions,''
Phys. Lett. B \textbf{759}, 336 (2016),
arXiv:1602.06953.

\bibitem{Boer2016}
D.~Boer, S.~Cotogno, T.~van Daal, P.~J.~Mulders, A.~Signori, and
Y.-J.~Zhou,
``Gluon and Wilson loop TMDs for hadrons of spin \(\leq1\),''
JHEP \textbf{10}, 013 (2016),
arXiv:1607.01654.

\bibitem{Gardi2013}
E.~Gardi, J.~M.~Smillie, and C.~D.~White,
``The non-Abelian exponentiation theorem for multiple Wilson lines,''
JHEP \textbf{06}, 088 (2013),
arXiv:1304.7040.

\bibitem{Blanes2009}
S.~Blanes, F.~Casas, J.~A.~Oteo, and J.~Ros,
``The Magnus expansion and some of its applications,''
Phys. Rept. \textbf{470}, 151 (2009),
arXiv:0810.5488.

\bibitem{EIS2012}
M.~G.~Echevarria, A.~Idilbi, and I.~Scimemi,
``Soft and collinear factorization and transverse momentum dependent parton
distribution functions,''
Phys. Lett. B \textbf{726}, 795 (2013),
arXiv:1211.1947.

\bibitem{Vladimirov2017}
A.~Vladimirov,
``Structure of rapidity divergences in soft factors,''
JHEP \textbf{04}, 045 (2018),
arXiv:1707.07606.

\bibitem{Cutkosky1960}
R.~E.~Cutkosky,
``Singularities and discontinuities of Feynman amplitudes,''
J. Math. Phys. \textbf{1}, 429 (1960).

\bibitem{Collins2011}
J.~Collins,
\emph{Foundations of Perturbative QCD},
Cambridge University Press (2011).

\bibitem{Lorce2011}
C.~Lorc\'e and B.~Pasquini,
``Quark Wigner distributions and orbital angular momentum,''
Phys. Rev. D \textbf{84}, 014015 (2011),
arXiv:1106.0139.

\bibitem{Marinho2008}
J.~A.~O.~Marinho, T.~Frederico, E.~Pace, G.~Salm\`e, and P.~U.~Sauer,
``Light-front Ward--Takahashi identity for two-fermion systems,''
Phys. Rev. D \textbf{77}, 116010 (2008),
arXiv:0805.0707.

\bibitem{Singh2010}
S.~Singh, R.~N.~C.~Pfeifer, and G.~Vidal,
``Tensor network decompositions in the presence of a global symmetry,''
Phys. Rev. A \textbf{82}, 050301(R) (2010),
arXiv:0907.2994.

\end{thebibliography}

\end{document}
